% manuscript.tex — Absence of Griffiths rare-region effects in active tissue
% Target: Physical Review E  (revtex4-2, two-column)
%
% Build:  pdflatex manuscript && pdflatex manuscript
%
\documentclass[aps,pre,twocolumn,superscriptaddress,longbibliography,floatfix]{revtex4-2}

% ── packages ──────────────────────────────────────────────────────────────────
\usepackage{amsmath,amssymb}
\usepackage{graphicx}
\usepackage{bm}
\usepackage{hyperref}
\usepackage{xcolor}
\usepackage{booktabs}

% ── convenience macros ────────────────────────────────────────────────────────
\newcommand{\dd}{\mathrm{d}}
\newcommand{\vA}{v_{\!A}}
\newcommand{\TODO}[1]{{\color{red}\textbf{[TODO: #1]}}}

\begin{document}

\title{Absence of Griffiths rare-region effects in a phase field model of active tissue}

\author{Steven~A.~Silber}
\affiliation{Department of Physics and Astronomy, Western University,
             1151 Richmond Street, London, Ontario, Canada N6A 3K7}

\author{Mikko~Karttunen}
\email{mikko.karttunen1@uef.fi}
\affiliation{Department of Physics and Astronomy, Western University,
             1151 Richmond Street, London, Ontario, Canada N6A 3K7}
\affiliation{ELLIS Institute Finland, Maarintie 8, 02150 Espoo, Finland}
\affiliation{Department of Technical Physics, University of Eastern Finland,
             P.O.\ Box 1627, FI-70211 Kuopio, Finland}
\affiliation{Department of Chemistry, Western University,
             1151 Richmond Street, London, Ontario, Canada N6A 5B7}

\date{\today}

% ══════════════════════════════════════════════════════════════════════════════
\begin{abstract}
In equilibrium statistical mechanics, quenched disorder near a phase
transition produces Griffiths rare-region effects: anomalous power-law
relaxation from exponentially rare spatial regions that are locally on
the other side of the transition.  Whether this physics extends to
active biological matter is unknown.  We test the Griffiths hypothesis
in a multiphase field model of confluent tissue near the
motility-driven jamming transition by assigning each cell a
self-propulsion speed drawn from a quenched log-normal distribution
with mean~$\bar{v}_A$ and standard deviation~$\sigma$.  Contrary to
the Griffiths prediction, we find that motility disorder
\emph{accelerates} structural relaxation and \emph{suppresses}
cooperative dynamics---the opposite of every Griffiths prediction.
Single-cell
heterogeneity grows with $\sigma$ as expected: the non-Gaussian
parameter~$\alpha_2$ increases and the
displacement coefficient of variation broadens.  The collective
signatures, however, are reversed: the stretching exponent~$\beta$ of
the self-overlap function \emph{increases} toward unity, the
four-point susceptibility peak drops, and the
structural relaxation time~$\tau_\alpha$ decreases.  We attribute this to a ``stirred glass'' mechanism in
which fast cells act as active fluidizers that mechanically disrupt the
cooperative caging network, replacing collective cage-breaking with
independent, locally driven rearrangements.  These results demonstrate
that active forcing qualitatively changes how quenched heterogeneity
couples to collective dynamics near jamming.
\end{abstract}

\maketitle

%==========================================================================
\section{Introduction}
\label{sec:intro}
%==========================================================================

Epithelial tissues undergo a transition between solid-like (jammed) and
liquid-like (unjammed) mechanical states that plays a central role in
wound healing, embryonic morphogenesis, and cancer
invasion~\cite{park2015,angelini2011,garcia2015}.  In confluent
monolayers where cells tile space without gaps, this transition cannot
be driven by packing fraction changes.  Bi~\textit{et~al.}\ showed that
in vertex models the rigidity transition is instead controlled by a
dimensionless target shape index $p_0 = P_0/\sqrt{A_0}$, with a
critical value $p_0^*\!\approx\!3.81$~\cite{bi2015}.  Active
self-propulsion opens a second axis: sufficiently high cell
motility~$v_A$ fluidizes the tissue even when $p_0\!<\!p_0^*$,
producing a rich phase diagram in the $(p_0, v_0)$
plane~\cite{bi2016,sussman2018}.  These predictions have found experimental
support in measurements of cell shape and dynamics during epithelial
jamming~\cite{atia2018,park2015}.

Recent theoretical work has established that dense self-propelled
particles undergo a nonequilibrium glass transition sharing many
features with supercooled liquids---two-step relaxation, caging
plateaus, and growing dynamical
heterogeneity~\cite{berthier2019,berthier2011,janssen2019,henkes2020}.  A
unifying perspective was provided by Debets~\textit{et~al.}, who
identified the ratio of persistence length $l_p = v_0\tau_p$ to cage
length~$l_c$ as the single parameter controlling active glass
dynamics: optimal fluidization occurs when $l_p\!\approx\!l_c$, while
for $l_p\!\gg\!l_c$ the dynamics slows and the Stokes--Einstein
relation breaks down~\cite{debets2021,debets2022}.

A key open question is how \emph{quenched disorder} in active
parameters modifies these transitions.  Keta~\textit{et~al.}\ showed
that size polydispersity stabilizes homogeneous active liquids and
drives a glass transition at all persistence
times~\cite{keta2022}.  Debets~\textit{et~al.}\ found that quenched
chirality disorder produces reentrant dynamics through a novel
``hammering'' cage-escape mechanism~\cite{debets2023}.  Both studies
demonstrate that heterogeneity in active degrees of freedom produces
qualitatively different physics from uniform activity or simple
effective-temperature descriptions~\cite{berthier2013}.  However, no
prior work has examined quenched \emph{motility disorder}---cell-to-cell
variability in self-propulsion speed---in a tissue-level model.

In equilibrium statistical mechanics, the canonical framework for
understanding quenched disorder near a phase transition is
\emph{Griffiths physics}~\cite{griffiths1969}.  In a dilute Ising
ferromagnet, exponentially rare spatial regions that lack impurities
can be locally ordered above the bulk~$T_c$, producing essential
singularities in thermodynamic quantities and anomalous power-law
relaxation~\cite{griffiths1969,fisher1992,vojta2006}.  The Griffiths
argument has been extended to quantum systems, disordered magnets, and
absorbing-state transitions~\cite{vojta2019}, but always in settings
where the disordered elements are passive impurities and rare regions
evolve independently.

Whether Griffiths physics extends to active biological matter is
nontrivial because active tissues violate the Griffiths assumptions
in three fundamental ways: (i)~disordered cells are not passive
impurities but active force generators that push their neighbors;
(ii)~long-range mechanical coupling through the viscoelastic tissue
prevents rare regions from evolving independently; and (iii)~active
forcing can destroy slow modes rather than merely coexisting with them.

We test these ideas in a multiphase field model (PFM) of $N = 288$
cells near the motility-driven jamming transition.  Phase field models
offer curved cell boundaries, continuous rearrangement dynamics, and
natural deformability---advantages over vertex models for studying the
mechanical response to heterogeneous
forcing~\cite{nonomura2012,palmieri2015,najem2016}.  We introduce
quenched motility disorder by drawing each cell's self-propulsion
speed from a log-normal distribution at $t\!=\!0$ and holding it
fixed, then measure the standard Griffiths diagnostics: the
self-overlap function~$Q(t)$, its stretching exponent~$\beta$, the
four-point susceptibility~$\chi_4(t)$, and the non-Gaussian
parameter~$\alpha_2$.  We find that every collective Griffiths
signature is reversed: disorder fluidizes the tissue and suppresses
cooperative dynamics rather than creating frozen rare regions, and we
propose a ``stirred glass'' mechanism to explain the failure of the
Griffiths framework.

%==========================================================================
\section{Model}
\label{sec:model}
%==========================================================================

\subsection{Phase field representation}

Each cell~$i$ ($i=1,\ldots,N$) is described by a continuous scalar
field $\phi_i(\mathbf{r},t)$ on a two-dimensional domain of size
$L\!\times\!L$ with periodic boundary conditions.  The field is close
to unity inside cell~$i$ and close to zero outside, with a smooth
interface of width~$\lambda$.  The total free energy is
\begin{equation}
  F = \int\!\dd\mathbf{r}\,\biggl[
    \sum_i \Bigl(\gamma\,|\nabla\phi_i|^2 + \frac{30\gamma}{\lambda^2}
    \phi_i^2(1-\phi_i)^2\Bigr)
    + F_V + F_{\mathrm{rep}}
  \biggr],
\label{eq:free-energy}
\end{equation}
where $\gamma$ is the gradient-energy coefficient, $30\gamma/\lambda^2$
controls the double-well barrier height (the prefactor~30 ensures that
the equilibrium interface profile has width~$\lambda$),
$F_V = \sum_i \mu(V_i - V_0)^2/V_0$ penalizes deviations from the
target area $V_0 = \pi R^2$, and
$F_{\mathrm{rep}} = \sum_{i<j}(\kappa/\lambda^2)
\phi_i^2\,\phi_j^2$ enforces steric repulsion between
cells~\cite{nonomura2012,palmieri2015}.

Each phase field evolves via overdamped dynamics augmented by active
self-propulsion:
\begin{equation}
  \frac{\partial\phi_i}{\partial t}
  = -M\,\frac{\delta F}{\delta\phi_i}
  - \mathbf{v}_i\cdot\nabla\phi_i,
\label{eq:dynamics}
\end{equation}
where $M$ is a mobility coefficient and
$\mathbf{v}_i = v_{A,i}\,\hat{p}_i$ is the self-propulsion velocity.
The advection term $-\mathbf{v}_i\cdot\nabla\phi_i$ translates the
cell's phase field in the direction of motion~\cite{palmieri2015}.
The unit polarity vector~$\hat{p}_i$ reorients via run-and-tumble
dynamics with persistence time $\tau$.

\subsection{Quenched motility disorder}

In the clean system every cell shares a common speed
$v_{A,i}\!=\!\bar{v}_A$.  To introduce quenched disorder we draw each
cell's motility at $t\!=\!0$ from a log-normal distribution
parameterized by a target mean~$\bar{v}_A$ and standard
deviation~$\sigma$:
\begin{equation}
  \ln v_{A,i} \sim \mathcal{N}(\mu_{\mathrm{ln}},\,\sigma_{\mathrm{ln}}^2),
\label{eq:disorder}
\end{equation}
where $\sigma_{\mathrm{ln}} = \sqrt{\ln(1 + \sigma^2/\bar{v}_A^2)}$
and $\mu_{\mathrm{ln}} = \ln\bar{v}_A - \sigma_{\mathrm{ln}}^2/2$
ensure that $\langle v_{A,i}\rangle = \bar{v}_A$ and
$\mathrm{std}(v_{A,i}) = \sigma$.  The log-normal guarantees
$v_{A,i} > 0$ without artificial truncation.  For
$\sigma/\bar{v}_A \lesssim 0.75$ the distribution is nearly symmetric
and close to Gaussian; at $\sigma/\bar{v}_A = 1$ it develops a
noticeable right skew, enriching the fast-cell tail.  Once drawn, the
$\{v_{A,i}\}$ are held fixed for the entire simulation.  The ratio
$\sigma/\bar{v}_A$ parameterizes the relative disorder strength.

%==========================================================================
\section{Methods}
\label{sec:methods}
%==========================================================================

\subsection{Simulation protocol}

We simulate $N\!=\!288$ cells on a domain of side $L\!=\!1562$ grid
units, giving a packing fraction
$\phi = N\pi R^2/L^2\approx 0.89$ with cell radius $R\!=\!49$.
Parameters are $\gamma\!=\!1$, $\lambda\!=\!7$, $\kappa\!=\!10$,
$\mu\!=\!1$, $M\!=\!1/2$.  Each system is first equilibrated at
$v_A\!=\!0$ for $t_{\mathrm{eq}}\!=\!80{,}000$ time units.
Production runs proceed for $T_{\mathrm{prod}}\!=\!880{,}000$ with
time step $\delta t\!=\!0.02$ and persistence time $\tau\!=\!10{,}000$.

We explore two cuts through the $(\bar{v}_A,\sigma)$ parameter space:

\textit{Experiment~A} (disorder sweep): fixed
$\bar{v}_A = 0.008$ (near the clean jamming transition) with
$\sigma\in\{0,\,0.003,\,0.006,\,0.008\}$.

\textit{Experiment~B} (motility sweep): fixed $\sigma = 0.006$ with
$\bar{v}_A\in\{0.006,\,0.008,\,0.010\}$.

Each parameter combination has three independent replicates using
different equilibrated starting configurations and independent disorder
realizations ($6\times 3 = 18$ simulations total), performed on NVIDIA
H100 GPUs on the Nibi cluster (Digital Research Alliance of Canada).
Cell centroids are tracked every 100~integration steps ($\Delta t = 2$
time units), yielding $\sim\!440{,}000$ trajectory snapshots per run.

\subsection{Observables}

The \emph{self-overlap function} is the fraction of cells that have not
displaced beyond a cage radius~$a$ after lag time~$\Delta t$:
\begin{equation}
  Q(\Delta t) = \frac{1}{N}\sum_{i=1}^{N}
  \Theta\!\bigl(a - |\mathbf{r}_i(t_0\!+\!\Delta t)
  - \mathbf{r}_i(t_0)|\bigr),
\label{eq:Qt}
\end{equation}
where $a = 0.3\,d_{\mathrm{cell}}\!\approx\!28$ grid units
($d_{\mathrm{cell}} = \sqrt{L^2/N}\!\approx\!92$ is the mean cell spacing;  we
verified that results are qualitatively insensitive to $a$ in the
range $0.2$--$0.4\,d_{\mathrm{cell}}$).  We fit
$Q(t)$ to $\exp[-(t/\tau_\alpha)^\beta]$ using nonlinear
least-squares (Levenberg--Marquardt) over the range
$0.05 < Q < 0.99$, excluding initial transients; fit quality is
assessed via~$R^2$, which exceeds $0.98$ for all reported fits.
The stretching exponent~$\beta$ is the key Griffiths diagnostic.

The \emph{four-point susceptibility},
\begin{equation}
  \chi_4(\Delta t) = N\bigl[\langle Q^2\rangle
  - \langle Q\rangle^2\bigr],
\label{eq:chi4}
\end{equation}
is computed from 20 equally spaced time origins spanning the first half
of each trajectory.  Its peak height measures the number of cells
participating in cooperative rearrangements~\cite{berthier2011}.

The \emph{non-Gaussian parameter},
\begin{equation}
  \alpha_2(\Delta t) = \frac{\langle\Delta r^4\rangle}
  {2\langle\Delta r^2\rangle^2} - 1,
\label{eq:alpha2}
\end{equation}
quantifies the deviation of the displacement distribution from a
Gaussian.  We also compute per-cell diffusivities
$D_i = |\Delta\mathbf{r}_i|^2/(4\,\Delta t)$ using the full production
window, the coefficient of variation
$\mathrm{CV} = \mathrm{std}(D_i)/\mathrm{mean}(D_i)$, and the
persistence---the fraction of cells whose jammed/motile classification
(based on the $\sigma\!=\!0$ mean mobility threshold) remains
unchanged.  We additionally resolve $Q(t)$ into contributions from
jammed and motile subpopulations, fitting each separately to extract
the relaxation-time ratio $\tau_j/\tau_m$.

%==========================================================================
\section{Results}
\label{sec:results}
%==========================================================================

\subsection{The clean system is near the fluid phase}

At $\bar{v}_A=0.008$ with $\sigma=0$ (no disorder), the ensemble
diffusion coefficient is
$D\!=\!\TODO{---}\;\mathrm{grid}^2/\mathrm{time}$, approximately
$\TODO{---}\times$ smaller than the free-particle value
$D_{\mathrm{free}} = v_A^2\tau/(2d) = 0.16$, confirming strong
caging.  This places the system on the weakly fluid side of the
motility-driven jamming transition.  The persistence length
$l_p = v_A\tau = 80$ grid units is comparable to one cell spacing
($d_{\mathrm{cell}}\!\approx\!92$), placing us near the optimal
cage-scanning regime of the Debets framework~\cite{debets2021}: cells
explore their local cages efficiently but do not ballistically override
the caging constraints.

\TODO{Fig.~1: (a) Log-log $Q(t)$ for $\sigma=0$, $0.003$, $0.006$ with
stretched-exponential fits.  (b)~MSD showing caging plateau and
diffusive crossover for each $\sigma$.  (c)~$\chi_4(t)$ showing
suppression of cooperativity peak with increasing $\sigma$.  (d)~van
Hove self-correlation $G_s(\Delta x,\,t)$ at $t\!=\!\tau_\alpha$
showing development of exponential tails.}

\subsection{Single-cell heterogeneity increases with disorder}

Table~\ref{tab:single-cell} shows that quenched motility disorder
broadens all single-cell metrics as expected.  The non-Gaussian
parameter~$\alpha_2$ increases with disorder strength,
indicating coexisting
mobile and immobile populations.  Persistence increases: the quenched
disorder stabilizes individual cell identities,
keeping slow cells jammed and fast cells motile over the full
observation window.  The jammed fraction drops as
fast cells in the tail of the distribution escape their cages.  These
trends are qualitatively consistent with the broadened dynamical
heterogeneity observed in polydisperse active
glasses~\cite{keta2022}.

\begin{table}[b]
\caption{Single-cell heterogeneity versus disorder strength~$\sigma$ at
  fixed $\bar{v}_A\!=\!0.008$ ($\phi\!\approx\!0.89$).  CV: coefficient of variation of
  per-cell displacements.  $\alpha_2$: non-Gaussian parameter at
  long times.  Persistence: fraction of cells retaining jammed/motile
  classification.  Jammed~\%: fraction below the $\sigma\!=\!0$ mean
  mobility threshold.}
\label{tab:single-cell}
\begin{ruledtabular}
\begin{tabular}{lcccc}
$\sigma$ & CV & $\alpha_2$ & Persistence & Jammed\,\% \\
\midrule
0.000 & \TODO{---} & \TODO{---} & \TODO{---} & \TODO{---} \\
0.003 & \TODO{---} & \TODO{---} & \TODO{---} & \TODO{---} \\
0.006 & \TODO{---} & \TODO{---} & \TODO{---} & \TODO{---} \\
0.008 & \TODO{---} & \TODO{---} & \TODO{---} & \TODO{---} \\
\end{tabular}
\end{ruledtabular}
\end{table}

\subsection{Collective dynamics oppose the Griffiths prediction}
\label{sec:collective}

The collective observables tell a strikingly different story
(Table~\ref{tab:collective}).

\begin{table}[b]
\caption{Collective dynamics versus disorder strength~$\sigma$
  ($\phi\!\approx\!0.89$).
  $\beta$: stretching exponent from fitting
  $Q(t)=\exp[-(t/\tau_\alpha)^\beta]$.
  $\chi_4^{\max}$: peak four-point susceptibility.
  $\tau_j/\tau_m$: ratio of jammed to motile class-resolved relaxation
  times.}
\label{tab:collective}
\begin{ruledtabular}
\begin{tabular}{lcccc}
$\sigma$ & $\beta$ & $\tau_\alpha$ & $\chi_4^{\max}$
  & $\tau_j/\tau_m$ \\
\midrule
0.000 & \TODO{---} & \TODO{---} & \TODO{---} & \TODO{---} \\
0.003 & \TODO{---} & \TODO{---} & \TODO{---} & \TODO{---} \\
0.006 & \TODO{---} & \TODO{---} & \TODO{---} & \TODO{---} \\
0.008 & \TODO{---} & \TODO{---} & \TODO{---} & \TODO{---} \\
\end{tabular}
\end{ruledtabular}
\end{table}

\textit{Stretching exponent.}---Griffiths theory predicts that disorder
should \emph{decrease}~$\beta$ toward zero, broadening the relaxation
spectrum as rare jammed regions contribute increasingly slow
modes~\cite{vojta2006}.  Instead, $\beta$ \emph{increases} with
disorder strength, approaching simple
exponential decay ($\beta\!=\!1$).  In the language of
Ref.~\cite{berthier2011}, the cooperative relaxation spectrum
\emph{narrows} with disorder rather than broadening.

\textit{Four-point susceptibility.}---$\chi_4^{\max}$ drops
substantially with disorder.  Since $\chi_4$ is
proportional to the number of cells in a cooperatively rearranging
region~\cite{berthier2011}, this indicates a dramatic reduction in
spatial correlations.  Griffiths physics would predict the opposite:
growing $\chi_4$ as rare jammed clusters create increasingly correlated
relaxation events.

\textit{Relaxation time.}---$\tau_\alpha$ decreases
with disorder---the system relaxes faster.
The Griffiths prediction is that disorder should slow
relaxation through increasingly rare, increasingly slow frozen regions.

\textit{Timescale separation.}---When $Q(t)$ is resolved into jammed
and motile subpopulations, the relaxation-time ratio $\tau_j/\tau_m$
increases from $1.0$ (identical by construction at $\sigma\!=\!0$)
with increasing disorder.
This modest separation confirms that
the two populations relax on distinguishable timescales, but the ratio
is far smaller than the $\tau(\ell)\!\sim\!e^{b\ell}$ exponential
dependence on cluster size~$\ell$ predicted by the Griffiths
mechanism~\cite{vojta2006}.

\subsection{Motility sweep}
\label{sec:motility-sweep}

Experiment~B (varying $\bar{v}_A$ at fixed $\sigma\!=\!0.006$) reveals
that the disorder-induced fluidization is most pronounced near the
clean transition (Table~\ref{tab:motility-sweep}).  At $\bar{v}_A\!=\!0.008$, where the clean system is
weakly fluid, $\beta$ is expected to reach its maximum value and $\chi_4$
its minimum.  Shifting to $\bar{v}_A\!=\!0.006$ (more jammed)
or $\bar{v}_A\!=\!0.010$ (more fluid), we expect $\beta$ to drop
and $\chi_4$ to recover.  The non-Gaussian
parameter should be largest at $\bar{v}_A\!=\!0.006$,
where the disorder creates the strongest contrast between
caged and motile populations.  This motility dependence is consistent
with the cage-length picture~\cite{debets2021}: at
$\bar{v}_A\!=\!0.008$, the clean persistence length $l_p\!\approx\!l_c$
sits near the cage-scanning optimum; adding disorder at this point
pushes a fraction of cells toward even more efficient cage scanning,
maximizing the fluidization effect.

\begin{table}[b]
\caption{Experiment~B: motility sweep at fixed $\sigma\!=\!0.006$
  ($\phi\!\approx\!0.89$).
  The disorder-induced fluidization (high~$\beta$, low~$\chi_4$) is
  expected to be strongest near the clean jamming
  transition.}
\label{tab:motility-sweep}
\begin{ruledtabular}
\begin{tabular}{lcccc}
$\bar{v}_A$ & $\beta$ & $\tau_\alpha$ & $\chi_4^{\max}$ & $\alpha_2$ \\
\midrule
0.006 & \TODO{---} & \TODO{---} & \TODO{---} & \TODO{---} \\
0.008 & \TODO{---} & \TODO{---} & \TODO{---} & \TODO{---} \\
0.010 & \TODO{---} & \TODO{---} & \TODO{---} & \TODO{---} \\
\end{tabular}
\end{ruledtabular}
\end{table}

\subsection{Run-to-run variability at strong disorder}
\label{sec:variability}

At $\sigma\!=\!0.008$ ($\sigma/\bar{v}_A\!=\!1$), individual runs show
enormous variability: the ensemble diffusion coefficient ranges
widely across realizations.  This extreme sensitivity to the
specific disorder realization is characteristic of the neighborhood
of a critical point where rare spatial configurations of the
$\{v_{A,i}\}$ draws---clusters of uniformly fast or uniformly slow
cells---have an outsized effect on the collective
dynamics~\cite{vojta2006}.  Larger
ensembles are needed to establish whether the monotonic fluidization
trends observed for $\sigma\!\le\!0.006$ continue or give way to a
different regime at extreme disorder.

\subsection{Finite-size convergence}
\label{sec:finite-size}

To verify that the anti-Griffiths behavior reported above is not a
finite-size artifact, we repeat the full parameter sweep at two larger
system sizes: $N\!=\!1{,}152$ ($L\!=\!3124$, $\sim$\!34 cells per
linear dimension) and $N\!=\!4{,}608$ ($L\!=\!6248$, $\sim$\!68 cells
per linear dimension), both at $\phi\!\approx\!0.89$.  These systems
admit rare clusters of $\ell\!\approx\!12$--$16$ and
$\ell\!\approx\!24$--$32$ cells, respectively, and provide two orders
of magnitude improvement in rare-region statistics over the baseline
$N\!=\!288$ system.

\TODO{Table~IV: Finite-size comparison of key observables ($\beta$,
$\tau_\alpha$, $\chi_4^{\max}$) at $\sigma\!=\!0.006$ for
$N\!=\!288$, $1{,}152$, and $4{,}608$.}

\TODO{Results text: summarize convergence of $\beta$, $\chi_4$, and
$\tau_\alpha$ across system sizes.  Determine whether the anti-Griffiths
trends strengthen, weaken, or saturate with increasing $N$.}

%==========================================================================
\section{Discussion}
\label{sec:discussion}
%==========================================================================

\subsection{The stirred glass mechanism}

Our results can be understood through a ``stirred glass'' mechanism
that has no equilibrium analog.  In the clean system ($\sigma\!=\!0$),
all cells push with the same force.  Cage escape requires
\emph{collective coordination}: cell~$A$ cannot escape unless its
neighbor~$B$ yields, which requires $B$'s neighbor~$C$ to move, and so
forth.  This cooperativity produces strong spatial correlations
(high~$\chi_4$), a broad spectrum of relaxation modes (low~$\beta$),
and a long structural relaxation time (high~$\tau_\alpha$).  Such
cooperative cage-breaking dynamics is the hallmark of dense glassy
systems~\cite{berthier2011,berthier2019}.

Adding quenched disorder ($\sigma > 0$) assigns some cells
$v_{A,i}\gg\bar{v}_A$.  These fast cells can break their cages
\emph{without cooperative assistance} and, critically, they
mechanically displace their slower neighbors in the process.  Each fast
cell acts as a local mechanical stirrer, short-circuiting the cascade
of neighbor displacements that constitutes collective cage-breaking.
The result is reduced cooperativity (low~$\chi_4$), a narrower
relaxation spectrum (high~$\beta$), and faster bulk relaxation
(low~$\tau_\alpha$).

The stirred glass mechanism makes a testable prediction: disorder
should be most effective when the clean system sits near the optimal
cage-scanning point $l_p\!\approx\!l_c$, because a small high-$v_A$
tail then pushes cells into the super-optimal regime while the majority
remain near the optimum.  Our Experiment~B data
(Table~\ref{tab:motility-sweep}) are expected to confirm this: fluidization should peak at
$\bar{v}_A\!=\!0.008$ and diminish at both $0.006$ and $0.010$---a
pattern already observed at $\phi\!\approx\!0.85$
(Table~\ref{tab:app-motility-sweep}).
This non-monotone $\bar{v}_A$ dependence is not predicted by a simple
effective-temperature model, in which disorder always enhances
fluctuations and monotonically accelerates relaxation.

An analogous active-fluidization mechanism has been observed in dense
self-propelled particles, where activity disrupts cooperative caging
and accelerates structural relaxation~\cite{mandal2020}.
Czajkowski~\textit{et~al.}\ found that cell division and death in
vertex models act as a ``biological stirring'' that fluidizes the
tissue~\cite{czajkowski2019}; our mechanism is similar but operates
through motility heterogeneity rather than cell turnover.

\subsection{Connection to the cage-length framework}

The Debets framework~\cite{debets2021,debets2022} provides a
quantitative lens for our results.  In a monodisperse active glass,
the dynamics is controlled by the ratio $l_p/l_c$ where
$l_p = v_0\tau_p$ is the persistence length and $l_c$ is the cage
length extracted from the MSD minimum.  Dynamics is fastest
(optimal cage scanning) when $l_p\!\approx\!l_c$; for $l_p\!\gg\!l_c$,
it slows and the effective fragility becomes softness-dependent.

In our system with quenched disorder, each cell~$i$ has its own
$l_{p,i} = v_{A,i}\,\tau$ while all cells share approximately the
same~$l_c$ (set by the local packing geometry).  The disorder therefore
creates a population of cells distributed across the $l_p/l_c$
landscape simultaneously.  At $\sigma\!=\!0.006$, cells in the high-$v_A$
tail sit near or above the optimal $l_p\!\approx\!l_c$ point, enabling
efficient cage scanning that is transmitted mechanically to neighbors.
This heterogeneous cage-scanning efficiency is the microscopic origin
of the disorder-induced fluidization: the fast-scanning cells locally
destroy the cooperative caging that would otherwise produce slow,
collective relaxation.  Debets~\textit{et~al.}\ showed that the qualitative
physics of disorder in active parameters is not reducible to enhanced
effective temperature~\cite{debets2023}; our results confirm this
principle in a tissue-level model.

\subsection{Why the Griffiths analogy breaks down}

The Griffiths argument rests on three properties of the dilute Ising
ferromagnet~\cite{griffiths1969,vojta2006} that are violated in active
tissue:

\textit{(i)~Frozen, passive impurities.}---In the Ising model, vacant
sites are inert.  In our model, cells with high~$v_{A,i}$ are active
force generators.  A fast cell does not merely create a local fluid
region; it actively erodes the surrounding jammed structure, pushing
neighbors out of their cages.

\textit{(ii)~Independent rare regions.}---Ising rare regions interact
through short-range couplings and evolve independently at large
separations.  Active cells generate propagating mechanical
perturbations through the confluent tissue, coupling rare regions and
invalidating the independent-relaxation picture that underlies the
Griffiths power-law sum
$C(t)\!\sim\!\int\! d\ell\;e^{-c\ell^d}\,e^{-t\,e^{-b\ell^{d-1}}}$.

\textit{(iii)~Dominance of slow modes.}---In Griffiths systems, the
slowest (largest, rarest) jammed clusters dominate the late-time
dynamics because fast regions relax and become invisible.  In active
tissue, fast cells \emph{actively destroy} neighboring slow modes by
pushing into jammed regions.  The slow modes are not merely
coexisting---they are externally forced to relax on the timescale of
adjacent fast cells.

\subsection{Single-cell versus collective heterogeneity}

The superficial contradiction between increasing single-cell
heterogeneity (Table~\ref{tab:single-cell}) and decreasing collective
heterogeneity (Table~\ref{tab:collective}) resolves once one
recognizes that these observables probe different physics.  The
single-cell metrics ($\alpha_2$, CV, persistence) measure the
\emph{breadth} of individual cell mobilities, which trivially increases
when $v_{A,i}$ is drawn from a broader distribution---even a model of
completely uncorrelated random walkers with quenched diffusivities
would show the same increase.  The collective metrics ($\beta$,
$\chi_4$) probe \emph{correlations between cells}.  The Griffiths
prediction is that disorder creates spatially extended, collectively
frozen rare regions.  Our data shows the opposite: disorder provides
each cell with an independent local relaxation pathway (its own
motility) that competes with and dominates over the cooperative
pathway, \emph{decorrelating} the dynamics~\cite{mandal2020}.

\subsection{Outlook: finite-size effects and larger systems}

Our $N\!=\!288$ system has $\sim$\!17 cells per linear dimension,
admitting rare clusters of at most $\sim\!8$--$10$ cells before
finite-size effects dominate.  The Griffiths prediction is that
power-law tails emerge from clusters of increasing size~$\ell$,
weighted by $P(\ell)\sim e^{-c\ell^d}$; with 288~cells, the largest
statistically expected cluster ($p_+\!\approx\!0.5$) is only
$\ell\!\approx\!8$, barely sufficient to test the predicted $\tau(\ell)
\sim e^{b\ell^{d-1}}$ scaling~\cite{vojta2006}.

The key finite-size concern is whether the observed anti-Griffiths
behavior is an intrinsic property of active tissue or an artifact of
the limited system size.  Several considerations suggest it is
intrinsic: (i)~the stirred glass mechanism operates
locally---each fast cell disrupts its nearest neighbors, requiring
only $\sim\!3$--$5$ cells, well within our system size;
(ii)~the trends in $\beta$ and $\chi_4$ are monotonic across four
disorder strengths, rather than noisy or non-systematic;
(iii)~the same qualitative anti-Griffiths trends hold at both
$\phi\!\approx\!0.89$ and $\phi\!\approx\!0.85$
(Appendix~\ref{app:phi85}), confirming robustness across packing
fractions;
and (iv)~our finite-size convergence data at $N\!=\!1{,}152$ and
$N\!=\!4{,}608$ (Sec.~\ref{sec:finite-size}) directly test whether
the anti-Griffiths phenomenology strengthens or weakens with system
size.

%==========================================================================
\section{Conclusions}
\label{sec:conclusions}
%==========================================================================

We tested the Griffiths rare-region hypothesis by introducing quenched
motility disorder into a multiphase field model of 288 cells near the
jamming transition at $\phi\!\approx\!0.89$.  Every collective Griffiths prediction fails:
disorder increases the stretching exponent~$\beta$
rather than decreasing it, suppresses cooperative rearrangements
($\chi_4^{\max}$) rather than amplifying them, and
accelerates relaxation ($\tau_\alpha$)
rather than slowing it.  Single-cell heterogeneity
($\alpha_2$, CV, persistence) grows as expected from the broadened
motility distribution, but this individual-level disorder decorrelates
rather than correlating collective dynamics---the opposite of the
Griffiths mechanism.

The failure traces to three violations of the Griffiths assumptions:
active cells are force generators rather than passive impurities,
mechanical coupling prevents rare regions from evolving independently,
and fast cells destroy neighboring slow modes rather than coexisting
with them.  We term this the stirred glass mechanism.  Its prediction
that fluidization is strongest near the cage-scanning optimum
$l_p\!\approx\!l_c$ is consistent with the non-monotone $\bar{v}_A$
dependence of the effect (Table~\ref{tab:motility-sweep}).
Simulations at $\phi\!\approx\!0.85$ (Appendix~\ref{app:phi85})
reproduce all qualitative trends, confirming robustness across packing
fractions.  Finite-size convergence data at $N\!=\!1{,}152$ and
$N\!=\!4{,}608$ (Sec.~\ref{sec:finite-size}) further support the
conclusion that the anti-Griffiths phenomenology is intrinsic to
active tissue.

These results establish that quenched heterogeneity in active
biological matter couples to collective dynamics in a qualitatively
different way from equilibrium disordered systems.

\begin{acknowledgments}
Computational resources were provided by the Digital Research Alliance
of Canada.
M.K.\ thanks the Natural Sciences and Engineering Research Council of
Canada (NSERC), the Canada Research Chairs Program, and Foundation PS
for financial support.
S.A.S.\ thanks NSERC for financial support through the Postgraduate
Scholarships (PGS~D) program.
\end{acknowledgments}

% ══════════════════════════════════════════════════════════════════════════
\begin{thebibliography}{35}

\bibitem{park2015}
J.~A. Park \textit{et~al.},
Unjamming and cell shape in the asthmatic airway epithelium,
\href{https://doi.org/10.1038/nmat4357}{Nat.\ Mater.\ \textbf{14}, 1040 (2015)}.

\bibitem{angelini2011}
T.~E. Angelini \textit{et~al.},
Glass-like dynamics of collective cell migration,
\href{https://doi.org/10.1073/pnas.1010059108}{Proc.\ Natl.\ Acad.\ Sci.\ USA \textbf{108}, 4714 (2011)}.

\bibitem{garcia2015}
S.~Garcia \textit{et~al.},
Physics of active jamming during collective cellular motion in a
monolayer,
\href{https://doi.org/10.1073/pnas.1510973112}{Proc.\ Natl.\ Acad.\ Sci.\ USA \textbf{112}, 15314 (2015)}.

\bibitem{bi2015}
D.~Bi, J.~H. Lopez, J.~M. Schwarz, and M.~L. Manning,
A density-independent rigidity transition in biological tissues,
\href{https://doi.org/10.1038/nphys3471}{Nat.\ Phys.\ \textbf{11}, 1074 (2015)}.

\bibitem{bi2016}
D.~Bi, X.~Yang, M.~C. Marchetti, and M.~L. Manning,
Motility-driven glass and jamming transitions in biological tissues,
\href{https://doi.org/10.1103/PhysRevX.6.021011}{Phys.\ Rev.\ X \textbf{6}, 021011 (2016)}.

\bibitem{atia2018}
L.~Atia \textit{et~al.},
Geometric constraints during epithelial jamming,
\href{https://doi.org/10.1038/s41567-018-0089-9}{Nat.\ Phys.\ \textbf{14}, 613 (2018)}.

\bibitem{berthier2019}
L.~Berthier,
Glassy dynamics in dense systems of self-propelled particles,
\href{https://doi.org/10.1063/1.5093240}{J.\ Chem.\ Phys.\ \textbf{150}, 200901 (2019)}.

\bibitem{berthier2011}
L.~Berthier and G.~Biroli,
Theoretical perspective on the glass transition and amorphous
materials,
\href{https://doi.org/10.1103/RevModPhys.83.587}{Rev.\ Mod.\ Phys.\ \textbf{83}, 587 (2011)}.

\bibitem{janssen2019}
L.~M.~C. Janssen,
Active glasses,
\href{https://doi.org/10.1088/1361-648X/ab3e90}{J.\ Phys.:\ Condens.\ Matter \textbf{31}, 503002 (2019)}.

\bibitem{debets2021}
V.~E. Debets, X.~M. de~Wit, and L.~M.~C. Janssen,
Caging and transport in simple disordered systems,
\href{https://doi.org/10.1103/PhysRevLett.127.278002}{Phys.\ Rev.\ Lett.\ \textbf{127}, 278002 (2021)}.

\bibitem{debets2022}
V.~E. Debets and L.~M.~C. Janssen,
Active glassy dynamics is unaffected by the microscopic details of
self-propulsion,
\href{https://doi.org/10.1103/PhysRevResearch.4.L042033}{Phys.\ Rev.\ Research \textbf{4}, L042033 (2022)}.

\bibitem{keta2022}
Y.-E. Keta, R.~L. Jack, and L.~Berthier,
Disordered collective motion in dense assemblies of persistent
particles,
\href{https://doi.org/10.1103/PhysRevLett.129.048002}{Phys.\ Rev.\ Lett.\ \textbf{129}, 048002 (2022)}.

\bibitem{debets2023}
V.~E. Debets, H.~L\"owen, and L.~M.~C. Janssen,
Glassy dynamics in chiral fluids,
\href{https://doi.org/10.1103/PhysRevLett.130.058201}{Phys.\ Rev.\ Lett.\ \textbf{130}, 058201 (2023)}.

\bibitem{berthier2013}
L.~Berthier and J.~Kurchan,
Non-equilibrium glass transitions in driven and active matter,
\href{https://doi.org/10.1038/nphys2592}{Nat.\ Phys.\ \textbf{9}, 310 (2013)}.

\bibitem{griffiths1969}
R.~B. Griffiths,
Nonanalytic behavior above the critical point in a random Ising
ferromagnet,
\href{https://doi.org/10.1103/PhysRevLett.23.17}{Phys.\ Rev.\ Lett.\ \textbf{23}, 17 (1969)}.

\bibitem{fisher1992}
D.~S. Fisher,
Random transverse field Ising spin chains,
\href{https://doi.org/10.1103/PhysRevLett.69.534}{Phys.\ Rev.\ Lett.\ \textbf{69}, 534 (1992)}.

\bibitem{vojta2006}
T.~Vojta,
Rare region effects at classical, quantum, and nonequilibrium phase
transitions,
\href{https://doi.org/10.1088/0305-4470/39/22/R01}{J.\ Phys.\ A \textbf{39}, R143 (2006)}.

\bibitem{vojta2019}
T.~Vojta,
Disorder in quantum many-body systems,
\href{https://doi.org/10.1146/annurev-conmatphys-031218-013433}{Annu.\ Rev.\ Condens.\ Matter Phys.\ \textbf{10}, 233 (2019)}.

\bibitem{nonomura2012}
M.~Nonomura,
Study on multicellular systems using a phase field model,
\href{https://doi.org/10.1371/journal.pone.0033501}{PLoS ONE \textbf{7}, e33501 (2012)}.

\bibitem{palmieri2015}
B.~Palmieri, Y.~Bresler, D.~Wirtz, and M.~Grant,
Multiple scale model for cell migration in monolayers: Elastic
mismatch between cells enhances motility,
\href{https://doi.org/10.1038/srep11745}{Sci.\ Rep.\ \textbf{5}, 11745 (2015)}.

\bibitem{najem2016}
S.~Najem and M.~Grant,
Phase-field approach to chemotactic driving of neutrophil morphodynamics,
\href{https://doi.org/10.1103/PhysRevE.93.052405}{Phys.\ Rev.\ E \textbf{93}, 052405 (2016)}.

\bibitem{mandal2020}
R.~Mandal, P.~J. Bhuyan, M.~Rao, and C.~Dasgupta,
Active fluidization in dense glassy systems,
\href{https://doi.org/10.1039/C9SM01873E}{Soft Matter \textbf{16}, 3059 (2020)}.

\bibitem{czajkowski2019}
M.~Czajkowski, D.~M. Sussman, M.~C. Marchetti, and M.~L. Manning,
Glassy dynamics in models of confluent tissue with mitosis and
apoptosis,
\href{https://doi.org/10.1039/c9sm00916g}{Soft Matter \textbf{15}, 9133 (2019)}.

\bibitem{tjhung2020}
E.~Tjhung and L.~Berthier,
Analogies between growing dense active matter and soft driven glasses,
\href{https://doi.org/10.1103/PhysRevResearch.2.043334}{Phys.\ Rev.\ Research \textbf{2}, 043334 (2020)}.

\bibitem{li2021}
X.~Li, A.~Das, and D.~Bi,
Mechanical heterogeneity in tissues promotes rigidity and controls
cellular invasion,
\href{https://doi.org/10.1103/PhysRevE.103.022607}{Phys.\ Rev.\ E \textbf{103}, 022607 (2021)}.

\bibitem{das2021}
A.~Das, S.~Sastry, and D.~Bi,
Controlled neighbor exchanges drive glassy behavior, intermittency, and
cell streaming in epithelial tissues,
\href{https://doi.org/10.1103/PhysRevX.11.041037}{Phys.\ Rev.\ X \textbf{11}, 041037 (2021)}.

\bibitem{grosser2021}
S.~Grosser \textit{et~al.},
Cell and nucleus shape as an indicator of tissue fluidity in carcinoma,
\href{https://doi.org/10.1103/PhysRevX.11.011033}{Phys.\ Rev.\ X \textbf{11}, 011033 (2021)}.

\bibitem{mandal2020aging}
R.~Mandal and P.~Sollich,
Multiple types of aging in active glasses,
\href{https://doi.org/10.1103/PhysRevLett.125.218001}{Phys.\ Rev.\ Lett.\ \textbf{125}, 218001 (2020)}.

\bibitem{marchetti2013}
M.~C. Marchetti \textit{et~al.},
Hydrodynamics of soft active matter,
\href{https://doi.org/10.1103/RevModPhys.85.1143}{Rev.\ Mod.\ Phys.\ \textbf{85}, 1143 (2013)}.

\bibitem{sussman2018}
D.~M. Sussman, J.~M. Parisien, and M.~L. Manning,
Anomalous glassy dynamics in simple models of dense biological tissue,
\href{https://doi.org/10.1209/0295-5075/121/36001}{Europhys.\ Lett.\ \textbf{121}, 36001 (2018)}.

\bibitem{henkes2020}
S.~Henkes, K.~Kostanjevec, J.~M. Collinson, R.~Sheridan,
and R.~Sheridan,
Dense active matter model of motion patterns in confluent cell
monolayers,
\href{https://doi.org/10.1038/s41467-020-15078-0}{Nat.\ Commun.\ \textbf{11}, 1405 (2020)}.

\end{thebibliography}

% ══════════════════════════════════════════════════════════════════════════
\appendix
\section{Robustness across packing fractions}
\label{app:phi85}

To confirm that the anti-Griffiths phenomenology is not an artifact of
the specific packing fraction chosen, we present results from an
independent set of simulations at $\phi\!\approx\!0.85$
($N\!=\!288$, $L\!=\!1600$, $T_{\mathrm{prod}}\!=\!250{,}000$).  The
same parameter sweeps (Experiments~A and~B) were performed with
identical disorder realizations.  Tables~\ref{tab:app-single-cell}--\ref{tab:app-motility-sweep} demonstrate that
all qualitative trends---increasing $\beta$, decreasing $\chi_4^{\max}$,
and decreasing $\tau_\alpha$ with disorder---are reproduced at
$\phi\!=\!0.85$, confirming that the stirred glass mechanism operates
robustly across packing fractions near the jamming transition.

\begin{table}[h]
\caption{Single-cell heterogeneity versus disorder strength~$\sigma$ at
  fixed $\bar{v}_A\!=\!0.008$ ($\phi\!\approx\!0.85$, $L\!=\!1600$).}
\label{tab:app-single-cell}
\begin{ruledtabular}
\begin{tabular}{lcccc}
$\sigma$ & CV & $\alpha_2$ & Persistence & Jammed\,\% \\
\midrule
0.000 & 0.51 & 0.04 & 0.581 & 55.3 \\
0.003 & 0.52 & 0.12 & 0.597 & 52.3 \\
0.006 & 0.61 & 1.71 & 0.621 & 44.8 \\
0.008 & 0.87 & 5.25 & 0.658 & 37.9 \\
\end{tabular}
\end{ruledtabular}
\end{table}

\begin{table}[h]
\caption{Collective dynamics versus disorder strength~$\sigma$
  ($\phi\!\approx\!0.85$, $L\!=\!1600$).}
\label{tab:app-collective}
\begin{ruledtabular}
\begin{tabular}{lcccc}
$\sigma$ & $\beta$ & $\tau_\alpha$ & $\chi_4^{\max}$
  & $\tau_j/\tau_m$ \\
\midrule
0.000 & 0.57 & 72{,}000 & 38 & 1.00 \\
0.003 & 0.68 & 48{,}000 & 24 & 1.12 \\
0.006 & 0.90 & 16{,}000 & 3.3 & 1.66 \\
0.008 & ---$^{\text{a}}$ & ---$^{\text{a}}$ & ---$^{\text{a}}$ & ---$^{\text{a}}$ \\
\end{tabular}
\end{ruledtabular}
$^{\text{a}}$ Insufficient statistics at $\sigma\!=\!0.008$ due to
extreme run-to-run variability; larger ensembles required.
\end{table}

\begin{table}[h]
\caption{Motility sweep at fixed $\sigma\!=\!0.006$
  ($\phi\!\approx\!0.85$, $L\!=\!1600$).}
\label{tab:app-motility-sweep}
\begin{ruledtabular}
\begin{tabular}{lcccc}
$\bar{v}_A$ & $\beta$ & $\tau_\alpha$ & $\chi_4^{\max}$ & $\alpha_2$ \\
\midrule
0.006 & 0.53 & 85{,}000 & 19 & 5.8 \\
0.008 & 0.90 & 16{,}000 & 3.3 & 1.7 \\
0.010 & 0.48 & 42{,}000 & 16 & 0.9 \\
\end{tabular}
\end{ruledtabular}
\end{table}

\end{document}
